\section{El Método Socrático}
El método socrático es una forma de diálogo cooperativo en la cual se busca estimular el pensamiento crítico a través de preguntas continuas. En lugar de proporcionar respuestas directas, el tutor socrático hace preguntas que guían al estudiante hacia la deducción de la solución por sí mismo.

\section{Inteligencia Artificial en la Educación}
Actualmente, existen herramientas como GitHub Copilot, diseñadas para mejorar la productividad de los desarrolladores profesionales. Sin embargo, estas no son adecuadas para un entorno estrictamente educativo inicial. Otros proyectos han intentado utilizar \textit{frameworks} conversacionales como ChatGPT para la enseñanza, pero habitualmente los alumnos caen en la tentación de pedir la respuesta final.

\section{Ingeniería de \textit{Prompts} (Prompt Engineering)}
Para lograr que un LLM se comporte de una forma específica, se emplean técnicas de \textit{prompt engineering}. Estas permiten inyectar directrices a nivel de sistema antes de iniciar la interacción con el usuario, configurando así la "personalidad" y las restricciones del modelo de IA subyacente.
